\begin{problem}[第37届全国中学生物理竞赛决赛][量子热机]{127}
    量子热机是利用量子物质作为工作物质进行循环的热机。下面以二能级原子系统为例描述量子热机的工作原理。二能级原子的平均能量定义为
    \begin{equation}
        \langle E \rangle = p _ {0} \cdot E _ {0} + p _ {1} \cdot E _ {1},
        \label{eq:1.p127}
    \end{equation}
    其中$E _ { 0 }$、$p _ { 0 }$和$E _ { 1 }$、$p _ { 1 }$分别表示原子处于基态和激发态的能量、概率。为简单起见，假设$E _ { 0 } = 0$，在循环过程中激发态与基态的能量差是一个可调参量。该原子处于能量为$E$的能态的概率满足玻尔兹曼分布
    \begin{equation}
        p \propto \mathrm{e} ^ {- E / (k _ {\mathrm{B}} T)},
        \label{eq:2.p127}
    \end{equation}
    其中$T$为热力学温度，$k _ { \mathrm { B } }$为玻尔兹曼常量。在准静态过程中，平均能量的变化为
    \begin{equation}
        \mathrm{d} \langle E \rangle = p _ {1} \mathrm{d} E _ {1} + E _ {1} \mathrm{d} p _ {1},
        \label{eq:3.p127}
    \end{equation}
    其中$p _ { 1 } \mathrm { d } E _ { 1 }$为能级变化引起的能量变化，对应外界对二能级原子系统所做的功$\mathrm{d}W$；$E _ { 1 } \mathrm { d } p _ { 1 }$为概率变化引起的能量变化，对应外界对二能级原子系统的传热$\mathrm{d}Q$。
    \begin{subproblem}
        将二能级原子系统与一个温度为$T$的热源接触，求热平衡时二能级原子系统处在基态的概率$p _ { 0 }$和激发态的概率$p _ { 1 }$。
    \end{subproblem}
    \begin{subproblem}
        经典奥托循环的$P$-$V$图如图2a所示，其中$A \to B$和$C \to D$是等容过程，$B \to C$和$D \to A$是绝热过程。试画出量子奥托循环过程中二能级的能量差$E _ { 1 }$与激发态概率$p _ { 1 }$的关系示意图，并计算量子奥托循环四个过程中的传热、内能增量和对外做功。
    \end{subproblem}
    \begin{subproblem}
        类似地，计算量子卡诺循环各个过程中的传热、内能增量和对外做功。假设量子卡诺循环中高温热源温度$T _ { h }$和低温热源温度$T _ { l }$分别与量子奥托循环的最高温度和最低温度相同，试比较量子奥托热机和量子卡诺热机的工作效率。
    \end{subproblem}
\end{problem}